\lecture{6}{2. Marts 2026}{Physical interpretation of the strain tensor}

\begin{problem}{3.2}
  The strain field at a point $P(x,y,z)$ in an elastic body is given by
  \begin{equation*}
    \Vec{\epsilon} = \begin{bmatrix}
      20 & 3 & 2\\
    3 & -10 & 5\\
    2 & 5 & -8\\
    \end{bmatrix} \times \num{e-6} 
  .\end{equation*}
  Determine the following values:
  \begin{enumerate}
    \item[(a)] The strain invariants
    \item[(b)] The principal strains
    \item[(e)] The mean normal strain and the strain deviator
  \end{enumerate}
\end{problem}

\begin{derivation}
  \begin{enumerate}
    \item[(a)]   The strain invariants are $Q_1, Q_2, Q_3$, which is the sum of the diagonal of $\Vec{\epsilon}$, the sum of the minors of the diagonal of $\Vec{\epsilon}$ and the determinant of $\Vec{\epsilon}$, respectively.
  The first strain invariant $Q_1$ is simply the sum of the diagonal entries of $\Vec{\epsilon}$ as
  \begin{equation*}
    Q_1 = 20 \cdot 10^{-6} + \left( -10 \cdot 10^{-6} \right) + \left( -8 \cdot 10^{-6} \right)  = \num{2e-6} 
  .\end{equation*}
  
  The second strain invariant $Q_2$ is the sum of the minors of the diagonal, as:
  \begin{equation*}
    Q_2 = \begin{vmatrix}
      20 \cdot 10^{-6} & 3 \cdot 10^{-6}\\
    3 \cdot 10^{-6}& -10 \cdot 10^{-6}\\
    \end{vmatrix} + \begin{vmatrix}
      20 \cdot 10^{-6}& 2\cdot 10^{-6}\\
    2\cdot 10^{-6} & -8\cdot 10^{-6}\\
    \end{vmatrix} + \begin{vmatrix}
      -10 \cdot 10^{-6} & 5 \cdot 10^{-6}\\
    5 \cdot 10^{-6}& -8 \cdot 10^{-6}\\
  \end{vmatrix} = -\num{3.18e-10}
.\end{equation*}
  
  The third strain invariant is the determinant of $\Vec{\epsilon}$, as
  \begin{equation*}
    Q_3 = \begin{vmatrix}
      20 \cdot 10^{-6} & 3 \cdot 10^{-6} & 2 \cdot 10^{-6}\\
    3 \cdot 10^{-6}& -10\cdot 10^{-6} & 5\cdot 10^{-6}\\
    2 \cdot 10^{-6}& 5 \cdot 10^{-6}& -8\cdot 10^{-6}\\
    \end{vmatrix} = \num{1.272e-15}
  .\end{equation*}
  
    \item[(b)] The principal strains are found by solving the eigenvalue problem:
      \begin{equation*}
\begin{bmatrix}
          20  & 3  & 2 \\
        3  & -10  & 5 \\
        2  & 5  & -8 \\
        \end{bmatrix} \rightarrow \{20.519,-14.1326,-4.38641\}
      .\end{equation*}
    \item[(c)] The mean normal strain is simply the mean of the three normal, i.e. $\epsilon_{ii}$ components of $\Vec{\epsilon}$:
      \begin{equation*}
        \epsilon_M = \frac{1}{3} \left( \epsilon_{11} + \epsilon_{22} + \epsilon_{33} \right) = \frac{1}{3} \cdot \left( 20 + \num{-10}  + \num{-8}  \right) \cdot \num{e-6}  = \num{0,666e-6} 
      .\end{equation*}
      Which means the mean normal strain tensor is:
      \begin{equation*}
        \Vec{\epsilon}_M = \begin{bmatrix}
          \num{0,666e-6}  & 0 & 0\\
        0 & \num{0,666e-6}  & 0\\
        0 & 0 & \num{0,666e-6} \\
        \end{bmatrix}
      .\end{equation*}
      The strain deviator tensor is then:
      \begin{equation*}
        \Vec{\epsilon}_D = \Vec{\epsilon} - \Vec{\epsilon}_M = \begin{bmatrix}
          20 - \num{0,666} & 3 & 2\\
        3 & -10 - \num{0,666}  & 5\\
        2 & 5 & -8 - \num{0,666} \\
        \end{bmatrix} \num{e-6} = \left(
\begin{array}{ccc}
 19.334 & 3. & 2. \\
 3. & -10.666 & 5. \\
 2. & 5. & -8.666 \\
\end{array}
\right) \num{e-6} 
      .\end{equation*}
  \end{enumerate}
\end{derivation}

\finalresult{
\begin{aligned}
&(\mathrm{a}) &  &\begin{aligned}
      Q_1 &= \num{2e-6}  \\
      Q_2 &= \num{-3,18e-10}  \\
      Q_3 &= \num{1,272e-15} 
    \end{aligned} \\ \\
&(\mathrm{b}) &  &\begin{aligned}
      \epsilon^{(1)} &= \num{20,519}  \\
      \epsilon^{(2)} &= \num{-14,132}  \\
      \epsilon^{(3)} &= \num{-4,386} 
    \end{aligned} \\ \\
&(\mathrm{c}) &  &\begin{aligned}
      \Vec{\epsilon}_M &= \begin{bmatrix}
        \num{0,666e-6}  & 0 & 0\\
      0 & \num{0,666e-6}  & 0\\
      0 & 0 & \num{0,666e-6} \\
      \end{bmatrix} \\
        \Vec{\epsilon}_D &= \begin{bmatrix}
          \num{19,334e-6} & \num{3e-6}  & \num{2e-6} \\
          \num{3e-6}  & \num{-10,666e-6}  & \num{5e-6} \\
          \num{2e-6}  & \num{5e-6}  & \num{-8,666e-6} \\
        \end{bmatrix}
    \end{aligned}
.\end{aligned}
}
\begin{problem}{3.3}
  With reference to Problem 3.2, consider a second coordinate system $x_i'$ in which $x_2'$ is parallel to $x_2$ and $x_1'$ and $x_3'$ arre defined by a \ang{60} counterclockwise rotation about $x_2$. Find the components of strain with reference to the $x_i'$ system using the strain field from Problem 3.2.
\end{problem}

\begin{derivation}
  We use the tensor transformation rule
\begin{equation*}
    \Vec{A}' = \Vec{R} \Vec{A} \Vec{R}^{\intercal}
  \end{equation*}
  where $\Vec{R}$ is the direction cosine matrix whose rows are the new basis vectors written in the old basis.
  In this case the tensor $\Vec{\epsilon}$ is to be rotated around the $x_2$-axis, which means:
  \begin{equation*}
    \Vec{R} = \begin{bmatrix}
      \cos \ang{60}  & 0 & \sin \ang{60} \\
    0 & 1 & 0\\
    - \sin \ang{60}  & 0 & \cos \ang{60} \\
    \end{bmatrix}
  .\end{equation*}
  So
  \begin{equation*}
    \Vec{\epsilon}' = \begin{bmatrix}
      \cos \ang{60}  & 0 & \sin \ang{60} \\
    0 & 1 & 0\\
    - \sin \ang{60}  & 0 & \cos \ang{60} \\
    \end{bmatrix} \cdot \begin{bmatrix}
      20 & 3 & 2\\
    3 & -10 & 5\\
    2 & 5 & -8\\
    \end{bmatrix} \cdot  \begin{bmatrix}
  \cos \ang{60}  & 0 & \sin \ang{60} \\
0 & 1 & 0\\
- \sin \ang{60}  & 0 & \cos \ang{60} \\
\end{bmatrix}^{\intercal} = \left(
\begin{array}{ccc}
  \num{5.86603} & \num{0}  & \num{5.13397} \\
 0 & -10 & 0 \\
 \num{-14.134} & 0 & \num{ -2.86603} \\
\end{array}
\right)
  .\end{equation*}
\end{derivation}

\finalresult{
\begin{bmatrix}
  \num{5,866} &0 & \num{5,134}  \\
0 & -10 & 0\\
\num{-14,134}  & 0 & \num{-2,866} \\
\end{bmatrix}}


\begin{problem}{3.6}
  Compute the infinitesimal strains, and sketch the deformed configuration of an initially rectangular volume element subject to the following displacement fields, where $A$, $B$, and $D_{ij} = \text{constants}$:
  \begin{enumerate}[label=(\alph*)]
    \item Simple extension: $u_1 = A_1 x_1; u_2 = u_3 = 0$.
    \item Simple shear: $u_1 = Bx_2; u_2 = u_3 = 0$.
    \item homogeneous deformation: $u_i = D_{ij} x_j$.
  \end{enumerate}
\end{problem}

\begin{derivation}
  \begin{enumerate}[label=(\alph*)]
    \item Here $u_1 = A_1 x_1$ and $u_2 = u_3 = 0$ leading to the following displacement gradient
      \begin{equation*}
        \nabla \Vec{u} = \begin{bmatrix}
          A_1 & 0 & 0\\
        0 & 0 & 0\\
        0 & 0 & 0\\
        \end{bmatrix}
      .\end{equation*}
      Here, the block simply stretches/contracts in the $x_1$-direction.
    \item Here $u_1 = B x_2$ and $u_2 = u_3 = 0$ leading to the following displacement gradient
      \begin{equation*}
        \nabla \Vec{u} = \begin{bmatrix}
          0 & B & 0\\
        0 & 0 & 0\\
        0 & 0 & 0\\
        \end{bmatrix}
      .\end{equation*}
      Which is symmetrized as:
      \begin{equation*}
        \Vec{\epsilon} = \frac{1}{2} \left( \nabla \Vec{u} + \left( \nabla \Vec{u} \right)^{\intercal} \right) = \frac{1}{2} \begin{bmatrix}
          0 & \frac{B}{2} & 0\\
        \frac{B}{2} & 0 & 0\\
        0 & 0 & 0\\
        \end{bmatrix}
      .\end{equation*}
      here we have a displacement in the $x_1$ direction proportional to $x_2$.
    \item Here we have the general case with deformation in all directions.
      The displacement is
      \begin{equation*}
        \Vec{\epsilon} = \frac{1}{2} \left( \Vec{D} + \left( \Vec{D} \right)^{\intercal} \right)
      .\end{equation*}
  \end{enumerate}
\end{derivation}

\finalresult{
\begin{aligned}
  &(\mathrm{a}) & \Vec{\epsilon} &= \begin{bmatrix}
          A_1 & 0 & 0\\
        0 & 0 & 0\\
        0 & 0 & 0\\
        \end{bmatrix} \\
  &(\mathrm{b}) &         \Vec{\epsilon} &= \frac{1}{2} \begin{bmatrix}
          0 & \frac{B}{2} & 0\\
        \frac{B}{2} & 0 & 0\\
        0 & 0 & 0\\
        \end{bmatrix} \\
  &(\mathrm{c}) &         \Vec{\epsilon} &= \frac{1}{2} \left( \Vec{D} + \left( \Vec{D} \right)^{\intercal} \right)
.\end{aligned}
}
